\input{sol_pre.tex}
\input{q_2023-02-19.tex}
\begin{document}
\input{sol_head.tex}

\Savol{1}
Yiğindi özgarmas bölsa, köpaytma sonlar bir-biriga eng yaqin bölganda eng katta: $55\cdot56=3080$, \ $3080-1=3079$.
\Javob{B}

\Savol{2}
Surat: $(1-3)+(5-7)+\dots+(2021-2023)$ — $506$ ta juft, har biri $-2$: \ $-1012$.\\
Maxraj: $(1-2)+\dots+(2021-2022)$ — $1011$ ta juft: \ $-1011$. \ Nisbat $\dfrac{1012}{1011}$.
\Javob{A}

\Savol{3}
\[\frac{10^{50}\left(10^5+10^4+1\right)}{10^{49}\left(1+10^4+10^5\right)}=10.\]
\Javob{C}

\Savol{4}
Birgalikda butun işni $6$ kunda bajaradi. $\dfrac1x+\dfrac1{x+5}=\dfrac16$ $\Rightarrow$ $x^2-7x-30=0$ $\Rightarrow$ $x=10$.
\Javob{A}

\Savol{5}
Dastlab temir $x$ kg. Ajratilgan qismda $25$ kg temir. Qolgan $300$ kg da $x-25$ kg:
\[\frac{x-25}{300}=\frac x{500}+0{,}2\ \Rightarrow\ 5x-125=3x+300\ \Rightarrow\ x=212{,}5.\]
\Javob{A}

\Savol{6}
$2022\cdot2024+1=2023^2$, \ $2024\cdot2030+9=2027^2$, \ $2026\cdot2028+1=2027^2$, \ $2020\cdot2026+9=2023^2$.\\
Köpaytma $1$: \ $\sqrt{1-1}+5=5$.
\Javob{D}

\Savol{7}
$a_5+a_{n-4}=a_1+a_n=24$; \ $S_n=\dfrac{24}2n=96$ $\Rightarrow$ $n=8$.
\Javob{D}

\Savol{8}
$b_3+b_4+b_5=q^2\left(b_1+b_2+b_3\right)$ $\Rightarrow$ $q^2=4$.\\
$q=2$: $7b_1=21$, $b_1=3<5$ — mos emas. \ $q=-2$: $3b_1=21$, $b_1=7$.
\[S_6=\frac{7\left(1-(-2)^6\right)}{1-(-2)}=\frac{7\cdot(-63)}3=-147.\]
\Javob{D}

\Savol{9}
Qavs içi, umumiy maxraj $ab(a-b)(a+b)$:
\[\frac{\left(a^2+b^2\right)\left(a^2-b^2\right)+b^3(a+b)-a^3(a-b)}{ab\left(a^2-b^2\right)}=\frac{ab\left(a^2+b^2\right)}{ab\left(a^2-b^2\right)}=\frac{a^2+b^2}{a^2-b^2}.\]
\[\frac{a^2+b^2}{a+b}\cdot\frac{(a-b)(a+b)}{a^2+b^2}=a-b.\]
\Javob{D}

\Savol{10}
$\dfrac1{1+2+\dots+k}=\dfrac2{k(k+1)}=2\left(\dfrac1k-\dfrac1{k+1}\right)$: \ $\displaystyle\sum_{k=1}^7=2\left(1-\frac18\right)=\frac74$.
\[\frac74x=7\ \Rightarrow\ x=4.\]
\Javob{B}

\Savol{11}
Şartlardan: $-1<a<0$, \ $b<-1$, \ $0<c<1$.\\
$ab>0$: \ $|b-c-1|=c+1-b$; \quad $bc<0$: \ $-|a+c+1|=-(a+c+1)$; \quad $ac<0$: \ $-|a+b+1|=a+b+1$.
\[c+1-b-a-c-1+a+b+1=1.\]
\Javob{A}

\Savol{12}
$\dfrac1{\sin106^\circ}+\dfrac{\cos106^\circ}{\sin106^\circ}=\dfrac{1+\cos106^\circ}{\sin106^\circ}=\dfrac{2\cos^253^\circ}{2\sin53^\circ\cos53^\circ}=\ctg53^\circ$.\\
$2\tg53^\circ\cdot\ctg53^\circ=2$.
\Javob{C}

\Savol{13}
\[\frac{(a-b)(a+b+3)}{(a+3)^2-b^2}=\frac{(a-b)(a+b+3)}{(a-b+3)(a+b+3)}=\frac{a-b}{a-b+3}=\frac47.\]
\Javob{D}

\Savol{14}
Çap tomon ösuvçi; $x=2$ da $4+9+16+25=54$. Yeçim $x>2$: \ $3,\dots,2023$ — $2021$ ta.
\Javob{A}

\Savol{15}
1) juft; \ 2) toq; \ 3) na juft, na toq; \ 4) $f(-x)=\log_2\frac1{x+\sqrt{1+x^2}}=-f(x)$ — toq. \ Jami $2$ ta.
\Javob{B}

\Savol{16}
$\left(x-\frac{2x}{5x+2}\right)^2+\frac{4x^2}{5x+2}=\frac95$; \ $t=x-\frac{2x}{5x+2}=\frac{5x^2}{5x+2}$, \ $\frac{4x^2}{5x+2}=\frac45t$:
\[t^2+\frac45t-\frac95=0\ \Rightarrow\ t=1\ \text{yoki}\ t=-\frac95.\]
$5x^2-5x-2=0$ ($D=65>0$) — $2$ ta; \ $5x^2+9x+3{,}6=0$ ($D=9>0$, $x=-0{,}6;\ -1{,}2$) — $2$ ta. Jami $4$.
\Javob{A}

\Savol{17}
$x+y=12$, $x-y=2$ $\Rightarrow$ $x=7$, $y=5$, \ $xy=35$.
\Javob{D}

\Savol{18}
$\pi^2<10$ $\Rightarrow$ maxraj manfiy: $\pi x>\sqrt{10}$, \ $x>\frac{\sqrt{10}}\pi\approx1{,}007$. \ $x=2,\dots,2023$ — $2022$ ta.
\Javob{D}

\Savol{19}
$1{,}025^n=9$ $\Rightarrow$ $n=\log_{1{,}025}9$.
\Javob{C}

\Savol{20}
Aniqlaniş: $4<x<5$; \ $\sqrt{5-x}>x-4$. Bu oraliqda butun son yöq.
\Javob{D}

\Savol{21}
$t=x^2+3x$: \ $(2x+3)\dfrac{t^2-16}t\ge0$, \ $t+4>0$ har doim:
\[\frac{(2x+3)(x+4)(x-1)}{x(x+3)}\ge0\ \Rightarrow\ x\in[-4;-3)\cup\left[-\tfrac32;0\right)\cup[1;\infty).\]
$[-4;5]$ dagi butunlar: $-4,\ -1,\ 1,2,3,4,5$: \ yiğindi $10$.
\Javob{D}

\Savol{22}
\[y'=-n\cos^{n-1}x\sin x\cos nx-n\cos^nx\sin nx=-n\cos^{n-1}x\,\sin(nx+x).\]
\Javob{C}

\Savol{23}
$\dfrac1{\cos kx\cos(k+1)x}=\dfrac{\tg(k+1)x-\tg kx}{\sin x}$. Yiğindi:
\[\frac{\tg2023x-\tg x}{\sin x}=\frac{\sin2022x}{\sin x\cos x\cos2023x}=\frac{2\sin2022x}{\sin2x\cos2023x}.\]
\Javob{C}

\Savol{24}
$AD$ — diametr, $AD=6$. \ $\angle ADC=\angle ABC=30^\circ$ (bir yoyga tiralgan), \ $\angle ACD=90^\circ$:
\[CD=AD\cos30^\circ=3\sqrt3.\]
\Javob{D}

\Savol{25}
$N$ — $AC$ örtasi: $S_{ADN}=\frac12S_{ADC}=\frac14S_{ABC}=12$. \ $GD=\frac13AD$ $\Rightarrow$ $S_{GDN}=\frac13\cdot12=4$.
\Javob{B}

\Savol{26}
Markaz $AD$ va $BC$ vatarlarning örta perpendikulyarlarida: $x=0{,}5$. Markaz $(0{,}5;k)$:
\[10{,}5^2+(k+3)^2=4{,}5^2+(5-k)^2\ \Rightarrow\ 16k=-74,\quad k=-4{,}625.\]
$R^2=10{,}5^2+1{,}625^2=\left(\frac{85}8\right)^2$, \ $R=10\frac58$.
\Javob{A}

\newpage
\Savol{27}
\fig{\begin{tikzpicture}[scale=0.6,font=\small]
\coordinate (L) at (-4,0); \coordinate (R) at (4,0); \coordinate (T1) at (-2,3.464); \coordinate (T2) at (2,3.464); \coordinate (B1) at (-2,-3.464); \coordinate (B2) at (2,-3.464);
\coordinate (M) at (0,3.464); \coordinate (P) at (-3,-1.732); \coordinate (Q) at (3,-1.732);
\coordinate (X1) at (-1,1.732); \coordinate (Y1) at (0.5,2.598); \coordinate (X2) at (1,-1.732); \coordinate (Y2) at (2.5,-0.866);
\fill[okgreen!30] (P)--(X1)--(Y1)--cycle; \fill[okgreen!30] (P)--(X2)--(Y2)--cycle;
\draw (L)--(T1)--(T2)--(R)--(B2)--(B1)--cycle; \draw (M)--(P)--(Q)--cycle; \draw (L)--(T2); \draw (B1)--(R); \draw (P)--(Y1); \draw (P)--(Y2);
\node[left] at (L) {$L$}; \node[above] at (T2) {$T$}; \node[below left] at (B1) {$B_1$}; \node[below right] at (B2) {$B_2$}; \node[right] at (R) {$R$}; \node[above] at (M) {$M$}; \node[left] at (P) {$P$}; \node[right] at (Q) {$Q$}; \node[above left] at (X1) {$X$}; \node[above] at (Y1) {$Y$};
\end{tikzpicture}}
$PQ=\frac{8+4}2=6$ ($LRB_2B_1$ trapetsiyaning örta çiziği), \ $S_{PMQ}=\frac{\sqrt3}4\cdot36=9\sqrt3$.\\
$\triangle MTY$: $\angle M=60^\circ$, $\angle T=30^\circ$ $\Rightarrow$ $\angle Y=90^\circ$, \ $MY=MT\cos60^\circ=1$.\\
$\triangle MTX$: $\angle M=120^\circ$, $\angle T=30^\circ$ $\Rightarrow$ teng yonli, $MX=MT=2$, \ $PX=4$.
\[S_{PXY}=\frac{PX}{PM}\cdot\frac{MY}{MQ}\cdot S_{PMQ}=\frac46\cdot\frac16\cdot9\sqrt3=\sqrt3.\]
Ikkinçi böyalgan uçburçak ham xuddi şunday: $\sqrt3$. \ Jami $2\sqrt3$.
\Javob{C}

\Savol{28}
$AA_1\perp A_1B_1$, $BB_1\perp A_1B_1$, ular orasidagi burçak $60^\circ$:
\[AB^2=AA_1^2+BB_1^2-2\cdot AA_1\cdot BB_1\cos60^\circ+A_1B_1^2=256+169-208+289=506.\]
\Javob{D}

\Savol{29}
$2a+b=0$; \ $f(2)-f'(2)=b+c=1$; \ $\frac a3+\frac b2+c=\frac23$.\\
$b=-2a$, $c=1+2a$: \ $\frac a3-a+1+2a=\frac23$ $\Rightarrow$ $a=-\frac14$, $b=\frac12$, $c=\frac12$. \ $b-a-c=\frac14$.
\Javob{A}

\Savol{30}
$\overrightarrow{BC}+\overrightarrow{CD}=\overrightarrow{BD}$, \ $\overrightarrow{D_1A_1}+\overrightarrow{A_1B}=\overrightarrow{D_1B}$. Qirra $1$: $\overrightarrow{BD}\cdot\overrightarrow{D_1B}=\overrightarrow{BD}\cdot\left(\overrightarrow{D_1D}+\overrightarrow{DB}\right)=0-2=-2$.
\[\cos\varphi=\frac{-2}{\sqrt2\cdot\sqrt3}=-\frac{\sqrt6}3.\]
\Javob{D}

\Savol{31}
$n(A\cup B)\ge9$ $\Rightarrow$ $n(A\cup B)'\le16-9=7$ ($B\subset A$ bölganda).
\Javob{A}

\Savol{32}
$10\cdot\mathrm C_9^2\cdot\mathrm C_7^2=10\cdot36\cdot21=7560$.
\Javob{B}

\newpage
\Savol{33}
\fig{\begin{tikzpicture}[scale=0.8,font=\small]
\coordinate (A) at (-2.83,0); \coordinate (B) at (2.83,0); \coordinate (C) at (0,2.83); \coordinate (C2) at (0,-2.83);
\fill[okgreen!30] (A)--(C)--(B)--(C2)--cycle; \draw (A)--(C)--(B)--(C2)--cycle;
\draw[ax] (-3.4,0)--(3.4,0); \draw[ax] (C)--(C2); \draw (0,0) circle (2); \draw[axs] (0,0)--(-1.414,1.414);
\node[left] at (A) {$A$}; \node[right] at (B) {$B$}; \node[above] at (C) {$C$}; \node[aux,below right] at (-0.9,1.0) {$\rho$};
\node[above left] at (-1.4,1.4) {$4$}; \node[aux,right] at (0,1.4) {$2\sqrt2$};
\end{tikzpicture}}
Bissektrisa ikkita teng yonli uçburçak hosil qiladi $\Rightarrow$ $\angle A=\angle B=45^\circ$, katetlar $4$. \ Aylaniş natijasida umumiy asosli ikki konus: $r=h=\frac{4\sqrt2}2=2\sqrt2$, yasovçi $l=4$.

\textbf{33.} $S=2\pi rl=2\cdot3\cdot2\sqrt2\cdot4=48\sqrt2$ — \textbf{D}.

\textbf{34.} $V=2\cdot\frac13\pi r^2h=\frac23\cdot3\cdot8\cdot2\sqrt2=32\sqrt2$ — \textbf{C}.

\textbf{35.} Içki şar radiusi — markazdan yasovçigaça masofa: $\rho=\frac{r\cdot h}l=\frac{8}4=2$. \ $V=\frac43\cdot3\cdot8=32$ — \textbf{B}.
\Javob{33 — D, \ 34 — C, \ 35 — B}

\newpage
\Savol{36}
$u=x^2-x-1$, \ $v=x^2-3x+2$, \ $u+v=2x^2-4x+1$. \ $u^3+v^3=(u+v)^3$ $\Leftrightarrow$ $3uv(u+v)=0$.\\
$u=0$: $x=\frac{1\pm\sqrt5}2$; \quad $v=0$: $x=1,\ 2$; \quad $u+v=0$: $x=1\pm\frac{\sqrt2}2$.

a) Barça $6$ ta ildiz turli.

b) Manfiy faqat $\frac{1-\sqrt5}2$; qolgan $5$ tasi musbat.
\Javob{a) 6; \quad b) 5}

\newpage
\Savol{37}
Maxraj: $\sin2^{x-2}\cos2^{x-2}=\frac12\sin2^{x-1}\ne0$. \ $\sin2^x=2\sin2^{x-1}\cos2^{x-1}$:
\[4\cos2^{x-1}=2\sqrt3\ \Rightarrow\ \cos t=\frac{\sqrt3}2,\qquad t=2^{x-1}\in\left(\tfrac12;16\right).\]
$t=\frac\pi6,\ \frac{11\pi}6,\ \frac{13\pi}6,\ \frac{23\pi}6,\ \frac{25\pi}6$ \ ($\frac\pi6\approx0{,}52>\frac12$, \ $\frac{35\pi}6>16$).

a) $5$ ta ildiz.

b) $x=1+\log_2t$: \ $x_{\max}-x_{\min}=\log_2\dfrac{25\pi/6}{\pi/6}=\log_225$.
\Javob{a) 5; \quad b) $\log_225$}

\newpage
\Savol{38}
$g'(x)=px^2+qx+r$: \ $4p+2q+r=7$, \ $16p+4q+r=3$, \ $36p+6q+r=-9$.\\
Ayirmalar: $12p+2q=-4$, $20p+2q=-12$ $\Rightarrow$ $p=-1$, $q=4$, $r=3$: \ $g'(x)=-x^2+4x+3$.

a) $k=g'(-2)=-4-8+3=-9$.

b) Statsionar nuqtalar: $x^2-4x-3=0$, Viyet böyiça yiğindi $4$.
\Javob{a) $-9$; \quad b) 4}

\newpage
\Savol{39}
\fig{\begin{tikzpicture}[scale=0.55,font=\small]
\coordinate (A) at (0,0); \coordinate (B) at (4,0); \coordinate (C) at (4,4); \coordinate (D) at (0,4); \coordinate (E) at (-1.6,5.6);
\fill[okgreen!30] (D)--(E)--(C)--cycle;
\draw (A)--(B)--(C)--(D)--cycle; \draw (B)--(E)--(C);
\draw (D)++(0.5,0) arc (0:135:0.5); \node[above right] at (0.25,4.35) {\scriptsize$135^\circ$};
\node[below left] at (A) {$A$}; \node[below right] at (B) {$B$}; \node[right] at (C) {$C$}; \node[below left] at (D) {$D$}; \node[above left] at (E) {$E$};
\end{tikzpicture}}
$\angle CDE=180^\circ-45^\circ=135^\circ$. \ $\triangle CDE$ da kosinuslar teoremasi ($CD=s$):
\[CE^2=s^2+9+2\cdot3s\cdot\frac{\sqrt2}2=s^2+3\sqrt2s+9.\]

a) $s=\frac{18}{\sqrt2}=9\sqrt2$: \ $CE^2=162+54+9=225$, \ $CE=15$.

b) $s^2+3\sqrt2s+9=29$ $\Rightarrow$ $s=\frac{-3\sqrt2+7\sqrt2}2=2\sqrt2$, \ $S=s^2=8$.
\Javob{a) 15; \quad b) 8}

\newpage
\Savol{40}
$(0;4)$: $m+d=4$; \ $(1;3)$: $n+4=3$, $n=-1$.

a) $n+m+d=3$.

b) $y=4-\sqrt[3]{x^2}=0$ $\Rightarrow$ $x=\pm8$. Simmetriya:
\[S=2\int_0^8\left(4-x^{\frac23}\right)dx=2\left(32-\frac35\cdot8^{\frac53}\right)=2\left(32-\frac{96}5\right)=\frac{128}5.\]
\Javob{a) 3; \quad b) $\dfrac{128}5$}

\newpage
\Savol{41}
\fig{\begin{tikzpicture}[scale=0.7,font=\small]
\coordinate (A) at (-5,0); \coordinate (C) at (3,0); \coordinate (H) at (0,0); \coordinate (B) at (0,5.196); \coordinate (D) at (1,3.464); \coordinate (E) at (-1.449,2.05); \coordinate (F) at (-0.639,3.025);
\fill[okgreen!30] (E)--(F)--(C)--cycle;
\draw (A)--(B)--(C)--cycle; \draw (B)--(H); \draw (A)--(D); \draw[axs] (B)--(E)--(C); \draw (E)--(F)--(C);
\draw (0,0.3)--(0.3,0.3)--(0.3,0);
\foreach \p in {D,E,F,H} \fill (\p) circle (2pt);
\node[left] at (A) {$A$}; \node[above] at (B) {$B$}; \node[right] at (C) {$C$}; \node[above right] at (D) {$D$}; \node[below left] at (E) {$E$}; \node[above left] at (F) {$F$}; \node[below] at (H) {$H$};
\end{tikzpicture}}
$\angle ADB=\angle AHB=90^\circ$ $\Rightarrow$ $A$, $B$, $D$, $H$ bir aylanada: \ $CD\cdot CB=CH\cdot CA=3\cdot8=24$.\\
$CB=CD+2$: $CD^2+2CD-24=0$ $\Rightarrow$ $CD=4$, $CB=6$.

a) $\triangle BEC$ töğri burçakli, $ED$ — balandlik: $CE^2=CD\cdot CB=24$, \ $CE=2\sqrt6$.

b) $\triangle EFC$ da $\angle F=90^\circ$: \ $FC=\sqrt{24-\left(12-6\sqrt3\right)}=\sqrt{12+6\sqrt3}$.
\[S=\frac12\sqrt{\left(12-6\sqrt3\right)\left(12+6\sqrt3\right)}=\frac12\sqrt{144-108}=3.\]
\Javob{a) $2\sqrt6$; \quad b) 3}

\newpage
\Savol{42}
\fig{\begin{tikzpicture}[scale=0.5,font=\small]
\coordinate (A) at (-3,0); \coordinate (B) at (1,-1.4); \coordinate (C) at (3,0.4); \coordinate (S) at (0.3,4.6); \coordinate (O) at (0.33,1.15);
\draw (A)--(B)--(C)--(S)--cycle; \draw (S)--(B); \draw[dashed] (A)--(C);
\fill[okgreen!30] (O) circle (1.0); \draw (O) circle (1.0); \draw[ax] (S)--(0.33,-0.33); \fill (O) circle (2pt);
\node[left] at (A) {$A$}; \node[below] at (B) {$B$}; \node[right] at (C) {$C$}; \node[above] at (S) {$S$}; \node[aux,right] at (O) {$r$};
\end{tikzpicture}}
$a=6\sqrt2$.

a) $V=\dfrac{a^3\sqrt2}{12}=\dfrac{432\sqrt2\cdot\sqrt2}{12}=72$.

b) Töla sirt $S_T=4\cdot\frac{\sqrt3}4a^2=72\sqrt3$. \ $r=\dfrac{3V}{S_T}=\dfrac{216}{72\sqrt3}=\sqrt3$:
\[S_{\text{şar}}=4\pi r^2=12\pi.\]
\Javob{a) 72; \quad b) $12\pi$}

\newpage
\Savol{43}
\fig{\begin{tikzpicture}[scale=0.7,font=\small]
\coordinate (A) at (0,0); \coordinate (B) at (1,0); \coordinate (C) at (3,3.464); \coordinate (D) at (2.5,4.33); \coordinate (E) at (0.5,4.33); \coordinate (F) at (-1,1.732);
\coordinate (P) at (-2,0); \coordinate (Q) at (5,0); \coordinate (R) at (1.5,6.062);
\fill[okgreen!30] (A)--(B)--(C)--(D)--(E)--(F)--cycle;
\draw[ax] (A)--(P)--(F); \draw[ax] (B)--(Q)--(C); \draw[ax] (D)--(R)--(E);
\draw (A)--(B)--(C)--(D)--(E)--(F)--cycle;
\node[below] at (A) {$A$}; \node[below] at (B) {$B$}; \node[right] at (C) {$C$}; \node[above right] at (D) {$D$}; \node[above left] at (E) {$E$}; \node[left] at (F) {$F$};
\node[aux,below] at (P) {$P$}; \node[aux,below] at (Q) {$Q$}; \node[aux,above] at (R) {$R$};
\end{tikzpicture}}
Barça burçaklar $120^\circ$. $AB$, $CD$, $EF$ davom ettirilsa, muntazam uçburçak $PQR$ hosil böladi; kesilgan uçburçaklar ham muntazam ($AF$, $BC$, $DE$ tomonli).\\
$PQ=FA+AB+BC=2+1+4=7$, \ $QR=BC+CD+DE=7$, \ $RP=DE+EF+FA=7$.

a) $DE=7-3-2=2$, \ $CD=7-4-2=1$.

b)
\[S=\frac{\sqrt3}4\left(7^2-2^2-4^2-2^2\right)=\frac{25\sqrt3}4.\]
\Javob{a) 1; \quad b) $\dfrac{25\sqrt3}4$}

\newpage
\Savol{44}
\fig{\begin{tikzpicture}[scale=0.38,font=\small]
\coordinate (O) at (0,0); \coordinate (A) at (8.485,8.485); \coordinate (K) at (5.796,10.039); \coordinate (L) at (10.039,5.796); \coordinate (M) at (3.107,11.593); \coordinate (N) at (11.593,3.107);
\draw[-{Stealth[length=4pt]}] (-1,0)--(15,0) node[right] {$x$}; \draw[-{Stealth[length=4pt]}] (0,-1)--(0,15) node[left] {$y$};
\draw (O)--(7.5,13); \draw (O)--(13,7.5); \draw (M)--(K)--(A)--(L)--(N); \draw[axs] (O)--(A); \draw[ax] (K)--(L); \draw[ax] (M)--(N);
\foreach \p in {A,K,L,M,N} \fill (\p) circle (5pt);
\node[below left] at (O) {$O$}; \node[right] at (A) {$A$}; \node[above right] at (K) {$K$}; \node[above right] at (L) {$L$}; \node[above left] at (M) {$M$}; \node[below right] at (N) {$N$}; \node[aux,above left] at (4.2,4.2) {$12$};
\end{tikzpicture}}
$OA=\sqrt{a^2+b^2}=12$. Çiziqlar $Ox$ bilan $60^\circ$ va $30^\circ$ hosil qiladi: $\angle KOL=30^\circ$.\\
$\angle OKA=\angle OLA=90^\circ$ $\Rightarrow$ $O$, $K$, $A$, $L$ — diametri $OA$ bölgan aylanada.

a) Sinuslar teoremasi: $KL=OA\sin\angle KOL=12\cdot\frac12=6$.

b) $K$, $L$ — $AM$ va $AN$ ning örtalari, $KL$ — $\triangle AMN$ ning örta çiziği: $MN=2KL=12$.
\Javob{a) 6; \quad b) 12}

\newpage
\Savol{45}
$a^2b=96$. Uy devoriga yopişgan devor $a\times b$; uning $4$ ta qirrasi ($2$ ta $b$, $2$ ta $a$) armaturasiz. Qolgan qirralar: $2$ ta $b$, $6$ ta $a$:
\[L=2b+6a=\frac{192}{a^2}+6a,\qquad L'=-\frac{384}{a^3}+6=0\ \Rightarrow\ a=4,\ b=6.\]

a) $L=12+24=36$ m.

b) Şişa: old devor $ab=24$ va ikki yon devor $2a^2=32$ — jami $56$ m$^2$, $560$\,\$. \ Tom: $ab=24$ m$^2$, $120$\,\$. \ Armatura $36$\,\$.
\[560+120+36=716\ \$.\]
\Javob{a) 36; \quad b) 716}

\ShowJavoblar
\end{document}
